% META: {"id": "TSPE-SUITLIM-EX-017", "chapitre": "TSPE-SUITES-LIMITES", "type_objet": "exercice", "capacites": ["TSPE-SUITLIM-C3"], "capacites_codes": ["C3"], "methodes": ["M3"], "parcours": 1, "competences": ["modeliser", "calculer"], "duree_min": 12, "mode_creation": "ex_nihilo", "sources_inspiration": [], "fichier_tex": "chapitres/TSPE-SUITES-LIMITES/exercices/TSPE-SUITLIM-EX-017.tex", "corrige_tex": "chapitres/TSPE-SUITES-LIMITES/corriges/TSPE-SUITLIM-CO-017.tex", "status": "generated"}
% BEGIN-VERIFY
% from sympy import *
% # Medicament: u_0=0, u_{n+1}=0.5*u_n+200
% # l = 200/(1-0.5) = 400
% # u_n = 400 - 400*0.5^n
% a = Rational(1, 2)
% b = 200
% l = Rational(400)
% assert l == b/(1-a)
% u0 = 0
% u1 = a*u0 + b
% assert u1 == 200
% u2 = a*u1 + b
% assert u2 == 300
% u3 = a*u2 + b
% assert u3 == 350
% END-VERIFY
\begin{exercice}{TSPE-SUITLIM-EX-017}{1}{12}
Un patient recoit chaque jour une dose de $200$~mg de medicament. Le corps elimine $50\,\%$ de la quantite presente chaque jour. On note $u_n$ la quantite de medicament (en mg) juste apres la prise du jour $n$, avec $u_0 = 0$ (pas de medicament avant le traitement, la premiere dose etant prise au jour $0$ : $u_0 = 200$).

On rectifie : $u_0 = 200$ (premiere prise) et $u_{n+1} = 0{,}5\,u_n + 200$.
\begin{enumerate}
  \item Calculer $u_1$, $u_2$ et $u_3$.
  \item Determiner la limite de $(u_n)$. Interpreter.
\end{enumerate}
\end{exercice}
